\documentclass[a4paper,12pt]{article}

\usepackage[T2A]{fontenc}
\usepackage[utf8]{inputenc}
\usepackage[russian]{babel}
\usepackage{amsmath}
\usepackage{array}

\begin{document}

\section*{6.2}

Та же схема, но:
\[
R \longleftrightarrow \text{квадратичный резонатор},
\qquad
U \text{ замкнуто}.
\]

\subsection*{1. При $R$}

Частота колебаний:
\[
f_0 = 1{,}014\ \text{МГц}.
\]

Напряжение:
\[
U_{\text{вых}} = 9658\ \text{мВ}.
\]

\subsection*{2. Воткнули квадратичный резонатор}

При
\[
C_2 = 20\ \text{пФ},
\qquad
A = 6000\ \text{мВ}.
\]

При
\[
C_2 = 47\ \text{пФ},
\qquad
A = 9130\ \text{мВ}.
\]

При этом подстроили контур.

\subsection*{3. Добротность}

\[
\Delta C = 20\ \text{пФ},
\qquad
f = 1{,}000\ \text{МГц}.
\]

\[
\Delta f = 994{,}0 - 974{,}0\ \text{кГц}
= 20\ \text{кГц}.
\]

\[
\Delta f_k = 1{,}001 - 0{,}995\ \text{МГц}
= 6\ \text{кГц}.
\]

\[
\frac{\Delta f_k}{\Delta f}
=
\frac{Q}{Q_k}.
\]

Отсюда:
\[
Q_k = Q \frac{\Delta f}{\Delta f_k}
\approx 83{,}3.
\]

\subsection*{4. При $C_s = 150\ \text{пФ}$}

\[
C_s = 150\ \text{пФ}.
\]

\[
\Delta f_k = 1{,}001 - 0{,}998\ \text{МГц}
= 3\ \text{кГц}.
\]

\[
\frac{\Delta f_k}{f_k}
=
\frac{C_k}{2C_s}.
\]

\[
C_k = 2C_s\frac{\Delta f_k}{f_k}
\approx 1\ \text{пФ}.
\]

\[
L_k =
\frac{1}{4\pi^2 f_k^2 C_k}
\approx 25\ \text{мГн}.
\]

\[
R_k =
\frac{1}{2\pi f_k C_k}
\approx 1{,}6 \cdot 10^5\ \Omega.
\]

\[
r_k =
\frac{R_k}{Q_k}
\approx 1{,}917\ \text{к}\Omega.
\]

\subsection*{5. Изменение частоты при $R = 300\ \Omega$}

Частота не меняется в пределах точности измерений при
\[
R = 300\ \Omega.
\]

\[
\begin{array}{|c|c|}
\hline
U_{\text{п}},\ \text{В} & f,\ \text{МГц} \\
\hline
-12 & - \\
\hline
-11 & - \\
\hline
-10 & 0{,}998 \\
\hline
-9 & 0{,}994 \\
\hline
-8 & - \\
\hline
-7{,}2 & 0{,}996 \\
\hline
\end{array}
\]

Вывод:
\[
f \approx \mathrm{const}.
\]

\section*{6.3}

\subsection*{1. Ёмкости}

\[
C_6 = 150\ \text{пФ},
\qquad
C_7 = 1{,}02\ \text{нФ},
\qquad
C_8 = 361\ \text{пФ}.
\]

Для последовательного соединения:
\[
C =
\left(
\frac{1}{C_6}
+
\frac{1}{C_7}
+
\frac{1}{C_8}
\right)^{-1}.
\]

По записи:
\[
C \approx 15 \cdot 10^{-12}\ \text{Ф}
=
15\ \text{пФ}.
\]

Частота генератора:
\[
f_g =
f_k\left(
1 + \frac{C_k}{2C}
\right)
=
1{,}034\ \text{МГц}.
\]

Вывод:
\[
\text{частота в пределах погрешности совпадает}.
\]

\subsection*{2. При изменении $C_7$}

\[
C_7^{*} = 471\ \text{пФ}.
\]

\[
\Delta f_g =
-
f_k
\frac{\Delta C_3 C_k}{2C_3^2}.
\]

Знак минус означает, что частота уменьшается.

\subsection*{3. Итог}

\[
\text{частота не меняется}.
\]

\end{document}
